%==================================================================

\subsection{Определение}
\leavevmode\par
\begin{defn}[слой, росток]
Пусть $\Shv{F}$ --- предпучок на $X$, $x\in X$. \emph{Слоем}
(пространством ростков над $x$) называется прямой предел
\begin{equation}\label{eq:stalk}
 \Shv{F}_x=\dirlim_{U\ni x}\Shv{F}(U),
\end{equation}
где $U$ пробегает открытые окрестности точки $x$, упорядоченные
включением (по контравариантности $\Shv{F}$ предел --- прямой над
категорией окрестностей).
\end{defn}

Элементы $\Shv{F}_x$ называются \emph{ростками} (зародышами). У
предельной конструкции есть наглядное описание, которое удобно
использовать на практике: элементы $\Shv{F}_x$ --- это классы
эквивалентности пар $(U,s)$, где $U\ni x$ открыто, $s\in\Shv{F}(U)$, а
эквивалентность
\[
 (U,s)\sim(V,t)
 \quad\Longleftrightarrow\quad
 \exists\, W\ni x,\ W\subset U\cap V:\ \res^{W}_{U}(s)=\res^{W}_{V}(t).
\]
Класс пары $(U,s)$ обозначают $s_x$ или $\tilde s(x)$: <<росток
сечения $s$ в точке $x$>>. Так записывается вычислительное правило:

\begin{quote}
\itshape Чтобы найти $s_x$, возьмите любую (достаточно малую)
окрестность $U$ точки $x$ и посмотрите, чем является $s$ на ней;
если на более малой окрестности $s$ совпал с другим сечением ---
ростки совпадают.
\end{quote}

\begin{exm}
Пусть $\Shv{F}(U)$ --- непрерывные $\mathbb{R}$-значные функции на
$U$, $\rho$ --- ограничение. Тогда $\Shv{F}_x$ --- это <<локальные
функции в точке $x$>>: два представителя $(U,f)$ и $(V,g)$ дают один
и тот же росток, если $f=g$ на некоторой окрестности $x$. Это
классическая конструкция зародыша функции из математического анализа.
\end{exm}

\begin{exm}[постоянный предпучок]
Если $\Shv{F}(U)=A$ для всех непустых $U$ и все $\rho$ ---
тождественные, то в пределе ничего не <<схлопывается>>:
$\Shv{F}_x=A$ для всякой $x$. Подробнее --- в
\S\,\ref{sec:const}.\end{exm}

\begin{risunok}{Как строится росток: вложенные окрестности
 $U\supset V\supset W\ni x$ и цепочка ограничений;
 росток --- класс эквивалентности <<стопки>> сечений.}
\begin{tikzpicture}[>=Stealth,font=\small]
  \draw[fig base] (0,0) ellipse (2.5 and 1.0);
  \draw[fig open] (0,0) ellipse (1.6 and 0.62);
  \draw[danger,line width=1pt,dotted] (0,0) ellipse (0.8 and 0.3);
  \fill (0,0) circle (2.8pt);
  \node[below right] at (0.12,-0.17) {$x$};
  \node[left] at (-2.5,0.42) {$U$};
  \node[left] at (-1.6,-0.45) {$V$};
  \node[above right] at (0.7,0.42) {$W$};
  % цепочка ограничений
  \node at (4.4,1.2)  {$\Shv{F}(U)$};
  \node at (4.4,0)    {$\Shv{F}(V)$};
  \node at (4.4,-1.2) {$\Shv{F}(W)$};
  \node at (6.7,-1.2) {$\Shv{F}_x$};
  \draw[->] (4.4,0.9) -- (4.4,0.3)
    node[midway,right,font=\footnotesize] {$\res$};
  \draw[->] (4.4,-0.9) -- (4.4,-0.3)
    node[midway,right,font=\footnotesize] {$\res$};
  \draw[->] (5.35,-1.2) -- (6.05,-1.2);
\end{tikzpicture}
\end{risunok}

Пояснение к рисунку. Слева --- вложенные окрестности точки $x$:
чёрная сплошная $U$, синяя пунктирная $V$, красная точками $W$;
в центре --- сама точка $x$. Справа --- что делает с этим
предпучок: сечение на $U$ ограничивается до $V$, затем до $W$
(стрелки $\rho$), и в пределе по такой цепочке получается росток
$\Shv{F}_x$. Росток и <<стопка>> сечений --- одно и то же:
два набора дают один и тот же росток, если они совпадают на
какой-то более мелкой окрестности.

\subsection{Почему слои важны}
Из лекции~1 мы помним главное правило: \emph{пучок определяется своими
слоями}, а точнее --- гомоморфизм пучков, который поэлементно (в каждом
слое) является изоморфизмом, есть изоморфизм. Это следствие того, что
аксиома (F1) позволяет <<восстановить>> глобальное сечение по его
росткам. Поэтому <<прокачивание>>, которое мы построим, ничего не
потеряет: слои $\Shv{F}_x$ и $\Shv{F}^{\#}_x$ будут совпадать.

\begin{bookbox}
\booksrc{Манин, \S\,2.1.6, с.~112 --- определение слоя
 $\Shv{P}_x=\dirlim \Shv{P}(U)$ и ростка.
\par\smallskip
 Годеман, гл.~II, \S\,1, п.~1.2, с.~130---133 --- накрывающее
 пространство, связанное с пучком (в оригинале \emph{espace étalé};
 переводчик предпочёл <<накрывающее пространство>>).
}
\end{bookbox}

%==================================================================
